\section{Conclusion and Future Work}
\label{sec:conclusion}

We have presented an automated, evidence-grounded pipeline that classifies
\textit{Defects4J} bugs into ODC Design/Code defect types using an LLM-driven
scientific classification approach: an enforced loop of hypothesis,
prediction, probe, and observation, inspired by scientific debugging and
adapted to classification rather than localization. The design's contributions are methodological as much
as empirical: a two-variable condition space that separates label space from
reasoning strategy, an open-set ODC formulation that turns taxonomy coverage into
a measured quantity, and a dual-mode pre-fix/post-fix evaluation that needs no
ground-truth labels. Together these address the five limitations of prior work:
full-dataset coverage through scalable automation, semantic ODC-type labels
rather than structural or patch-level ones, pre-fix operation without post-fix
dependency, controlled oracle leakage, and a self-consistent evaluation
framework.

We formalised five research questions (distribution, taxonomy coverage,
four-level accuracy, component contribution, and pre-fix/post-fix divergence)
and built a reproducible study toolchain that computes each from condition-tagged
artefacts over a shared evidence corpus. Development-validation runs confirm the
pipeline and its analysis layer produce correct, well-formed metrics end to end;
the quantitative findings are populated by the confirmatory run described in
Section~\ref{sec:setup}.

\smallskip
Future work spans five directions:
\begin{enumerate}
  \item \textbf{Full-dataset classification}: scale to the entire active
        \textit{Defects4J} set, producing the first complete ODC-labelled version
        of the benchmark.
  \item \textbf{Multi-model consensus}: classify with several LLMs and
        aggregate by majority vote to reduce single-provider variance
        \citep{colavito2024llm}.
  \item \textbf{Completing the opener/closer attribute set}: implement
        ODC's full field-defect procedure (Trigger from the report's usage
        scenario, Activity via the \S2.1 mapping) for the subset of bugs whose
        reports carry sufficient usage detail; recover Age via VCS archaeology
        that reconstructs when defective code was introduced; and validate the
        Impact attribute of Section~\ref{sec:impact} against trained human
        raters.
  \item \textbf{ODC-guided repair}: use classifications to select repair
        strategies, connecting classification output to automated repair tools
        (e.g.\ \textit{Checking} bugs to conditional-guard templates)
        \citep{durieux2018automatic}.
  \item \textbf{Cross-dataset and cross-language extension}: adapt the pipeline
        to other benchmarks and languages, reusing the language-independent
        classification logic.
\end{enumerate}

%% ─── Acknowledgements ────────────────────────────────────────────────────────
